%==============================================================================
%  feed-twin: governing equations, constitutive closure, and numerical method
%
%  REVISION 2 -- audited line-by-line against the implementation.
%
%  The mathematics below is what the code does, including the constants. The
%  prose is deliberately skeletal. Two kinds of comment mark where you write:
%
%    % EXPAND  -- narrative belongs here; the equations are already stated.
%    % NOTE:   -- what the paragraph needs to *say*, and sometimes why.
%
%  Compile: pdflatex twice (cleveref), or `tectonic feed-twin-physics.tex`.
%==============================================================================
\documentclass[11pt,a4paper]{article}

\usepackage[margin=1in]{geometry}
\usepackage{amsmath,amssymb,amsthm}
\usepackage{booktabs}
\usepackage{siunitx}
\usepackage{graphicx}
\usepackage{hyperref}
\usepackage{cleveref}
\usepackage{enumitem}
\usepackage{array}

\sisetup{detect-all, per-mode=symbol}
% Units siunitx does not ship. psi is the working unit on a stand, and the
% paper quotes it wherever the hardware does.
\DeclareSIUnit{\psi}{psi}
\DeclareSIUnit{\psig}{psig}
\DeclareSIUnit{\inch}{in}

\newtheorem{proposition}{Proposition}
\newtheorem{remark}{Remark}
\newtheorem{assumption}{Assumption}
\newtheorem{definition}{Definition}

\newcommand{\mdot}{\dot{m}}
\newcommand{\Qdot}{\dot{Q}}
\newcommand{\dd}[2]{\frac{\partial #1}{\partial #2}}
\newcommand{\eff}{\mathcal{E}}          % NTU effectiveness (not roughness)
\newcommand{\Nfix}{\mathcal{N}_{\mathrm{fix}}}
\newcommand{\Nfree}{\mathcal{N}_{\mathrm{free}}}
\newcommand{\code}[1]{\texttt{#1}}

\title{A Lumped-Parameter Digital Twin of a Pressure-Fed\\
       Rocket Propellant Feed System:\\
       Governing Equations, Constitutive Closure, and Numerical Method}
\author{CalSTAR Propulsion}
\date{\today}

\begin{document}
\maketitle

\begin{abstract}
% EXPAND
% NOTE: One paragraph. What is modelled (pressurant supply through regulator
% to tanks to injector to chamber), at what fidelity (lumped, quasi-steady
% lines, exact vessel thermodynamics), and the three things a reader should
% take away: (i) the algebraic network is closed exactly inside every
% derivative evaluation; (ii) every constitutive relation is adapted from a
% published source, not fitted; (iii) every optional model reproduces the
% baseline bit-for-bit when off.
We formulate a pressure-fed liquid propellant feed system as a semi-explicit
index-1 differential--algebraic system on a directed graph. Vessel inventories
evolve differentially in $(m, U)$; the hydraulic network is closed algebraically
at every right-hand-side evaluation by a row-scaled, damped Newton iteration
whose Jacobian is the graph incidence matrix plus one finite-difference scalar
per branch. Constitutive closure is by adapted, published correlations. We give
the complete equation set with the constants the implementation uses, the
thermal transport model, the engine coupling, both time-marching schemes (a
stiff implicit integrator for batch studies and an explicit, sub-stepped
scheme with an RC stability bound for the interactive cockpit), the
drawing-to-network construction rules, and the verification methodology.
Every closure's validity limit is stated where it is introduced and collected
in \cref{sec:limits}.
\end{abstract}

\tableofcontents

%==============================================================================
\section{Scope and modelling posture}
%==============================================================================
% EXPAND
% NOTE: Say what the tool is *for* -- pad-side and burn-time prediction of
% tank pressure, pressurant budget, mixture ratio drift, and chamber pressure
% for a small pressure-fed engine -- and what it is *not* for: injector
% internal flow, acoustics, manifold maldistribution, two-phase feed.

The model is \emph{lumped-parameter}: there is no mesh, no spatial
discretisation, and consequently no CFL condition. A feed system is a directed
graph whose edges carry one-dimensional, quasi-steady constitutive relations and
whose vertices carry either a prescribed pressure or a mass balance. The
algebraic state dimension is of order $10^1$--$10^2$; the differential state
dimension is of order $10^1$.

Three commitments shape everything below.

\begin{enumerate}[label=(C\arabic*)]
\item \textbf{Adapted, not invented.} Every correlation is taken from a
      published source through a thin adapter onto a reference implementation
      (\code{fluids}, \code{ht}, CoolProp). Where a constant cannot be adapted
      it is \emph{calibrated} at import against the reference implementation
      rather than transcribed (\cref{sec:choking}).
\item \textbf{Provenance is part of the state.} Every hardware number is a
      tuple $(\text{value}, \text{unit}, \text{source}, \text{reference})$ with
      $\text{source}\in\{\text{measured},\text{manufacturer},\text{estimated},\text{default}\}$,
      ranked in that order. There is no implicit source. Every result is
      stamped with the versions of every library that can move a number.
\item \textbf{Optional physics is inert by default.} A model that can be
      switched off must, when off, reproduce the previous answer bit-for-bit;
      and when on against a geometry that declares nothing for it, likewise
      (\cref{sec:toggle}).
\end{enumerate}

%==============================================================================
\section{Nomenclature}
%==============================================================================
\begin{center}
\small
\begin{tabular}{lll}
\toprule
Symbol & Meaning & Unit \\
\midrule
$\mathcal{N},\mathcal{B}$ & node and branch sets; $\Nfix,\Nfree$ fixed and free nodes & --- \\
$p_i$ & pressure at node $i$ (absolute) & \si{\pascal} \\
$T_i$ & temperature at node $i$ (given to the steady solve) & \si{\kelvin} \\
$\mdot_b$ & mass flow in branch $b$, positive upstream$\to$downstream & \si{\kilogram\per\second} \\
$d_i$ & external demand drawn from node $i$ & \si{\kilogram\per\second} \\
$\rho,\mu,k,c_p,c_v$ & density, viscosity, conductivity, heat capacities & SI \\
$h,u$ & specific enthalpy, specific internal energy & \si{\joule\per\kilogram} \\
$\gamma$ & heat-capacity ratio $c_p/c_v$ & --- \\
$R_s$ & specific gas constant $R/M$ & \si{\joule\per\kilogram\per\kelvin} \\
$D$ & bore (internal diameter) & \si{\metre} \\
$L$ & length & \si{\metre} \\
$\varepsilon$ & absolute roughness & \si{\metre} \\
$K$ & resistance coefficient, referred to a stated bore & --- \\
$f_D$ & Darcy friction factor & --- \\
$\mathrm{Re},\mathrm{Pr},\mathrm{Nu},\mathrm{Bi}$ & Reynolds, Prandtl, Nusselt, Biot & --- \\
$C_v,K_v$ & valve flow coefficients & --- \\
$\sigma$ & actuator travel fraction, $0$ shut, $1$ open & --- \\
$\chi(\sigma)$ & inherent valve characteristic & --- \\
$x, x_T, F_\gamma, Y$ & IEC 60534 pressure-drop ratio, its terminal factor, $\gamma/1.40$, expansion factor & --- \\
$C_d$ & discharge coefficient & --- \\
$m,U$ & vessel mass, total internal energy & \si{\kilogram}, \si{\joule} \\
$V_u,V_\ell$ & ullage and liquid volume & \si{\cubic\metre} \\
$A_i$ & gas--liquid interface area & \si{\square\metre} \\
$\Pi$ & expansion power $p\,\mathrm{d}V_u/\mathrm{d}t$ & \si{\watt} \\
$\beta,\alpha$ & thermal effusivity $\sqrt{k\rho c}$, diffusivity $k/(\rho c)$ & SI \\
$\eff$ & heat-exchanger effectiveness & --- \\
$c^*,C_F,A_t$ & characteristic velocity, thrust coefficient, throat area & SI \\
$\Lambda$ & injector stiffness $\Delta p_{\mathrm{inj}}/p_c$ & --- \\
$\kappa$ & Jacobian slope floor & \si{\pascal\second\per\kilogram} \\
$\mathcal{K}$ & choked-row stiffness & \si{\pascal\second\per\kilogram} \\
$\tau_{RC}$ & regulator--ullage time constant & \si{\second} \\
$\Gamma$ & explicit-coupling safety factor & --- \\
\bottomrule
\end{tabular}
\end{center}

%==============================================================================
\section{Network formulation}\label{sec:network}
%==============================================================================

\subsection{Graph and state}
Let $G=(\mathcal{N},\mathcal{B})$ be a directed graph. Each branch $b$ has an
upstream node $\mathrm{up}(b)$ and a downstream node $\mathrm{dn}(b)$. A node
is \emph{fixed} if it carries a prescribed pressure (vessel ullages and
outlets, ambient, the chamber) and \emph{free} otherwise. All pressures inside
the model are \textbf{absolute}; the stand's transducers cancel atmosphere, so
every pressure a person reads from or types into the twin is \textbf{gauge}
(psig, vented $=0$), converted at the API boundary and nowhere else.
Differential pressures are the same in either. Every node carries a
fluid species and a temperature; \textbf{the steady solve holds temperature
fixed} and solves for pressure only. Temperature is updated between solves by
\cref{sec:thermal}.

The algebraic unknown vector is
\begin{equation}
  \mathbf{y} = \bigl(\{p_i\}_{i\in\Nfree},\;\{\mdot_b\}_{b\in\mathcal{B}'}\bigr)
  \in\mathbb{R}^{n},
  \qquad n = |\Nfree'| + |\mathcal{B}'|,
\end{equation}
where the primes denote the sets after the reductions of \cref{sec:reduce}.

\subsection{Governing equations}
\begin{align}
  g_i(\mathbf{y}) &=
    \sum_{b:\,\mathrm{dn}(b)=i}\mdot_b - \sum_{b:\,\mathrm{up}(b)=i}\mdot_b - d_i = 0,
    && i\in\Nfree', \label{eq:mass}\\[2pt]
  g_b(\mathbf{y}) &=
    \bigl(p_{\mathrm{up}(b)}-p_{\mathrm{dn}(b)}\bigr)
    - \Delta p_b^{\mathrm{tot}}\bigl(\mdot_b;\,\boldsymbol{\phi}_b\bigr) = 0,
    && b\in\mathcal{B}', \label{eq:branch}
\end{align}
where $\boldsymbol{\phi}_b$ is the fluid state and signal set evaluated
\textbf{at the branch's upstream node}: $\rho,\mu,\gamma,R_s$ at
$(p_{\mathrm{up}}, T_{\mathrm{up}})$, together with $p_{\mathrm{sat}}(T_{\mathrm{up}})$
and $p_{\mathrm{crit}}$ (\cref{sec:phase}). The signed total is
\begin{equation}
  \Delta p_b^{\mathrm{tot}}(\mdot)
    = \operatorname{sgn}(\mdot)\,\Delta p_b(|\mdot|) + \Delta p_b^{\mathrm{static}},
  \label{eq:signed}
\end{equation}
for every component except the regulator (\cref{sec:regulator}).

\subsection{Reductions before the solve}\label{sec:reduce}
% NOTE: These three steps are why a stand with the whole panel shut is a
% zero-dimensional problem rather than a singular one. Say that most of a real
% feed system is stubs.
\begin{enumerate}[label=(R\arabic*)]
\item \textbf{Isolation.} A valve commanded to $\sigma\le\sigma_{\mathrm{shut}}=10^{-3}$
      is an open circuit: its branch is removed from $\mathcal{B}$ and its flow
      set to zero. A caller may add branches (a dry tank's outlets).
\item \textbf{Dead-end peeling.} Iteratively, any free node with zero demand
      and exactly one live branch is a stub: node and branch are removed. A free
      node left with \emph{no} live branch (cut off by an isolated valve) is a
      stub whose live end lies across the removed branch. Peeling repeats until
      no stub remains.
\item \textbf{Exact back-fill.} After the solve each stub carries $\mdot=0$ and
      $p_{\mathrm{stub}} = p_{\mathrm{live}} \mp \Delta p^{\mathrm{static}}$.
      A stub whose component holds a drop at zero flow ($\Delta p_b(0)>0$: a
      check valve below cracking, a regulator at lockup) is reported
      \emph{indeterminate}, and indeterminacy propagates along a chain of stubs.
\end{enumerate}

\subsection{Initial guess, warm start, and row scaling}
\begin{equation}
  p_i^{(0)} = \frac{1}{|\Nfix|}\sum_{k\in\Nfix}p_k,\qquad
  \mdot_b^{(0)} = \mdot_0 = 10^{-3}\ \si{\kilogram\per\second}
  \quad\text{(non-zero on purpose; see \cref{sec:regular})}.
\end{equation}
A caller may override any entry with a previous solution (warm start). The
residual is dimensionally mixed, so convergence is tested on row-scaled entries,
\begin{equation}
  s_i = \max\Bigl(\max_{j\in\Nfree'}|d_j|,\ \mdot_0,\ 10^{-3}\Bigr)
  \ \text{for mass rows},\qquad
  s_b = \max_{k\in\Nfix}|p_k|\ \text{for branch rows},
  \label{eq:scales}
\end{equation}
\begin{equation}
  \|\mathbf{g}\|_{\mathrm{sc}} = \max_{r\notin\mathcal{I}}\frac{|g_r|}{s_r} < \mathrm{tol},
  \qquad
  \mathrm{tol} = \begin{cases}10^{-6} & \text{standalone steady}\\ 10^{-5} & \text{inside a batch transient}\\ 10^{-4} & \text{live cockpit}\end{cases}
  \label{eq:norm}
\end{equation}
with $\mathcal{I}$ the exclusion set of \cref{sec:regular}. Row scaling does not
change the Newton step ($\mathbf{J}^{-1}\mathbf{g}$ is invariant); it changes
only the test.

\subsection{Jacobian structure}\label{sec:jac}
% EXPAND: this is why a step costs one derivative per component.
Mass rows are the (constant) incidence matrix:
$\partial g_i/\partial\mdot_b\in\{+1,-1,0\}$, $\partial g_i/\partial p_j = 0$.
An unchoked branch row has two constant entries and one scalar:
\begin{equation}
  \dd{g_b}{p_{\mathrm{up}}}=+1,\qquad
  \dd{g_b}{p_{\mathrm{dn}}}=-1,\qquad
  \dd{g_b}{\mdot_b}=-\,\dd{\Delta p_b^{\mathrm{tot}}}{\mdot_b}
  \approx -\frac{\Delta p_b^{\mathrm{tot}}(\mdot_b+\delta)-\Delta p_b^{\mathrm{tot}}(\mdot_b-\delta)}{2\delta},
  \label{eq:fd}
\end{equation}
a \textbf{central} difference with $\delta = \max(10^{-6}|\mdot_b|,\,10^{-9})$.
No correlation is differentiated analytically. If either evaluation raises,
the slope is set to $\kappa$ (\cref{sec:regular}). Residual and derivative are
assembled in one pass so each component is evaluated three times per
iteration, not six.

\subsection{Regularisation at zero flow}\label{sec:regular}
Quadratic losses give $\partial\Delta p/\partial\mdot\to 0$ as $\mdot\to0$, so
the branch row of \cref{eq:fd} is singular there. The assembled entry is
floored, sign-preserving:
\begin{equation}
  J_{bb} = -\operatorname{sgn}\!\Bigl(\dd{\Delta p_b}{\mdot_b}\Bigr)
           \max\Bigl(\Bigl|\dd{\Delta p_b}{\mdot_b}\Bigr|,\ \kappa\Bigr),
  \qquad \kappa = 10^{-3}\ \si{\pascal\second\per\kilogram}.
\end{equation}
Every unchoked branch with $|\partial\Delta p/\partial\mdot|<\kappa$ enters
$\mathcal{I}$: its row is \emph{unresolvable}, not merely small, and it is
excluded from \cref{eq:norm} and from the backtracking merit function. All such
branches are reported at convergence.

\begin{proposition}[Regularisation cannot mask a conservation error]
$\mathcal{I}$ contains only branch rows. Every mass balance \cref{eq:mass} is
tested, so conservation holds to $\mathrm{tol}$ at every node regardless of how
many branches are floored.
\end{proposition}

\subsection{Choked branch rows}\label{sec:chokedrow}
A branch is written in choked form if the component reports that the
available drop exceeds its critical value (\cref{sec:choking}; for a regulator,
that it has shut, with ceiling $0$). A second route exists for an iterate whose
\emph{flow} guess the component cannot price --- a gas orifice asked to pass
more than sonic --- and it is taken only when the gradient is forward,
$p_{\mathrm{up}}>p_{\mathrm{dn}}$; against the gradient the flow is priced just
under the ceiling and the row stays invertible so Newton can reverse it.
\begin{remark}
The forward-gradient condition is load-bearing. Without it a Newton iterate
that overshot a solenoid's sonic ceiling was pinned there regardless of the
pressures, and a pinned row does not depend on $p_{\mathrm{dn}}$: on a helium
press manifold this converged with sonic flow running \emph{up} a
\SI{16}{\psi} gradient into a tank already above its regulator's setpoint,
\SI{26}{\psi} per pair of sub-steps into a three-gram ullage, ratcheting the
tank to \SI{745}{\psig} against a \SI{550}{\psig} regulator.
\end{remark}
The choked row is
\begin{equation}
  g_b = \mathcal{K}\bigl(\mdot^{\mathrm{ceil}}_b - \mdot_b\bigr),\qquad
  \dd{g_b}{\mdot_b} = -\mathcal{K},\quad
  \dd{g_b}{p_{\mathrm{up}}} = \mathcal{K}\frac{\mdot^{\mathrm{ceil}}_b}{p_{\mathrm{up}}},\quad
  \dd{g_b}{p_{\mathrm{dn}}} = 0,
  \qquad \mathcal{K}=10^{8}\ \si{\pascal\second\per\kilogram},
  \label{eq:chokerow}
\end{equation}
the $p_{\mathrm{up}}$ derivative following from $\mdot^{\mathrm{ceil}}\propto p_{\mathrm{up}}$
for an ideal gas. A ceiling of exactly zero means \emph{shut} (a regulator that
has closed) and is honoured only when the first test fires.

\subsection{Damped Newton iteration}
Solve $\mathbf{J}\Delta\mathbf{y} = -\mathbf{g}$ by sparse LU (a failed
factorisation is a failed solve, reported as ``a node with no path to a fixed
pressure''). Then backtrack:
\begin{equation}
  \lambda = \max\Bigl\{2^{-j}: j=0,\dots,19,\ \
    \min_{i\in\Nfree'}\bigl(y_i + 2^{-j}\Delta y_i\bigr) > 0,\ \
    \|\mathbf{g}(\mathbf{y}+2^{-j}\Delta\mathbf{y})\|_{\mathrm{sc}} < \|\mathbf{g}(\mathbf{y})\|_{\mathrm{sc}}\Bigr\}.
\end{equation}
If no $j$ qualifies, a last-resort step $\lambda=0.01$ is taken with node
pressures clamped to $\max(0.01\,p_i,\ \SI{1}{\pascal})$; the iteration cap
(100 standalone, 30 live) is what ends a stalled solve.

\subsection{Physicality guard}
A converged residual on an iterate with any $p_i\le0$ is rejected as a failed
solve. Newton knows nothing of absolute zero, and a stiff branch can satisfy
\cref{eq:branch} at a large negative pressure that the property layer will
refuse several frames later.

%==============================================================================
\section{Fluid state at a node}\label{sec:phase}
%==============================================================================
% NOTE: The phase rule is the single largest hidden assumption in the steady
% solve. Say that phase is a *declaration* made by whatever built the network,
% and that the temperature rule is only the fallback for a node nobody
% declared.

Branch properties are evaluated from the upstream node's $(p,T)$ through a
property chain (\cref{sec:props}). Every node carries a declared phase
$\varphi\in\{\text{liquid},\text{gas},\varnothing\}$, set from topology by the
drawing reader (\cref{sec:pid}, D2). With the default \code{multiphase = False},
\begin{equation}
  (\rho,\mu,\gamma,p_{\mathrm{sat}}) = \begin{cases}
    \text{values at }(p,T), & \varphi=\text{gas},\\[2pt]
    \text{saturated-liquid values at } T\ \text{(independent of $p$)}, & \varphi=\text{liquid},\ \text{or } \varphi=\varnothing \text{ and } T<T_{\mathrm{crit}},\\[2pt]
    \text{values at }(p,T),\ p_{\mathrm{sat}}=0, & \varphi=\varnothing \text{ and } T\ge T_{\mathrm{crit}}.
  \end{cases}
  \label{eq:phaserule}
\end{equation}
The liquid branch exists so that a propellant leg can never be silently priced
as gas because a temperature reached it wrong (LOX at
\SI{1140}{\kilogram\per\cubic\metre} against gaseous oxygen at \num{40}), and
it makes a feed line's properties constant across a solve, which is cached
(quantised in $T$). The declaration exists because the fallback is by
\emph{temperature alone} and cannot tell a subcritical vapour from its liquid.

\begin{remark}[Consequences of the rule]\label{rem:phase}
(i) A liquid is priced saturated; real feed lines run subcooled and are a few
percent stiffer. (ii) An \emph{undeclared} node below $T_{\mathrm{crit}}$ is
priced as liquid whatever it holds --- an ethanol-vapour line
($T_{\mathrm{crit}}=\SI{514}{\kelvin}$) would come out as liquid ethanol ---
which is why every node the drawing can classify is declared, and only nodes
the drawing cannot place (a bare junction reached by nothing) fall through.
(iii) Pressurant lines are unaffected either way: N$_2$ and He are supercritical
at every pad temperature.
\end{remark}

%==============================================================================
\section{Constitutive closure: branch models}\label{sec:branches}
%==============================================================================

\subsection{Dynamic head and the resistance formalism}
\begin{equation}
  \Delta p = K\cdot\tfrac12\rho V^2,\qquad
  V=\frac{\mdot}{\rho\,\frac{\pi}{4}D^2},\qquad
  \mathrm{Re}=\frac{4|\mdot|}{\pi D\mu}.
  \label{eq:head}
\end{equation}
Every $K$ below is referred to a stated bore.

\subsection{Straight pipe (\code{darcy})}
\begin{equation}
  \Delta p_{\mathrm{pipe}} = \Bigl(f_D\frac{L}{D}+K_{\mathrm{minor}}\Bigr)\tfrac12\rho V^2,
  \qquad
  \Delta p^{\mathrm{static}} = \rho g\,\Delta z,
  \label{eq:darcy}
\end{equation}
\begin{equation}
  f_D = \begin{cases}
    0, & \mathrm{Re}\le0,\\
    64/\mathrm{Re}, & 0<\mathrm{Re}<\mathrm{Re}_{\mathrm{lam}}=2040,\\
    f_D^{\text{method}}(\mathrm{Re},\varepsilon/D), & \mathrm{Re}\ge\mathrm{Re}_{\mathrm{lam}},
  \end{cases}
\end{equation}
where the turbulent method is a named option (default \emph{Clamond}, an
essentially exact solution of Colebrook--White). The set of admissible names is
read from \code{fluids} at import so the two cannot disagree. Static head is
carried separately because it does not change sign with the flow
(\cref{eq:signed}).

\subsection{Discrete fitting (\code{fitting})}
\begin{equation}
  K_{\mathrm{fit}} = \mathcal{F}_{\mathrm{kind}}\bigl(D,\mathrm{Re},\varepsilon,f_D;\,\boldsymbol\theta\bigr),
\end{equation}
from a registry of adapters onto \code{fluids.fittings}: rounded bends at a
stated angle and $r/D$ (default $5$), sharp contraction and sudden expansion
(both need $D_2$), sharp entrance, normal exit, Crane converging tee run and
branch, Crane ball, gate, globe and swing-check valves, and Crane's
$K = n\,f_T$ forms ($n=30$ for a 90$^\circ$ elbow, $16$ for 45$^\circ$).
\begin{remark}[The fitting-$K$ trap]
A tabulated $K$ already prices the flow through the fitting's own body. The
friction term of \cref{eq:darcy} must therefore use tube length only; using an
end-to-end run length with the elbows also counted as $K$ pays for them twice.
\end{remark}

\subsection{Hardline bend (\code{bend})}
$K$ from \code{fluids.bend\_rounded} with the \emph{actual} centreline radius
$r_c$ and angle. No curved-friction term is added: the correlation already
contains it, and between $r_c/D=1$ and $5$ the loss changes severalfold.
$r_c<r_{\min}$ is reported as a violation, never enforced.

\subsection{Flexible hose (\code{flex\_hose})}
\begin{equation}
  \Delta p = \Bigl(f_D^{\mathrm{hose}}\frac{L}{D}+K_{\mathrm{end}}\Bigr)\tfrac12\rho V^2,
  \qquad
  f_D^{\mathrm{hose}} = \begin{cases}
    f_D^{\mathrm{curved}}\bigl(\mathrm{Re},D,D_c=2r_{\mathrm{inst}},\varepsilon\bigr)\cdot\zeta, & r_{\mathrm{inst}}>0,\\
    f_D(\mathrm{Re},\varepsilon/D)\cdot\zeta, & \text{straight},
  \end{cases}
\end{equation}
with $\zeta$ the convolution multiplier ($1$ smooth-bore; user-declared for
convoluted hose) and $K_{\mathrm{end}}$ the two crimped end fittings (default
$0.5$). Static and dynamic minimum bend radii are checked; the dynamic one is
a warning.

\subsection{Segmented line (\code{segmented}) and the authority ladder}\label{sec:segments}
\begin{equation}
  \Delta p_{\mathrm{line}} = \sum_{s=1}^{S}\Delta p_s + \sum_{s=1}^{S-1}\Delta p^{\mathrm{trans}}_{s\to s+1},
  \qquad
  \Delta p^{\mathrm{static}} = \rho g\sum_s\Delta z_s\ \ (\text{or the line total if no segment states one}).
\end{equation}
\begin{center}
\begin{tabular}{clc}
\toprule
Rank & Method & $\Delta p_s$ \\
\midrule
1 & \code{curve}       & linear interpolation of a measured $\Delta p(\mdot)$; \emph{refuses} outside its range \\
2 & \code{measured\_K} & $K_s\,\tfrac12\rho V_s^2$, $K_s$ from a flow bench \\
3 & \code{itemised}    & $\bigl(f_D L_s^{\mathrm{tube}}/D_s+\sum_j K_j\bigr)\tfrac12\rho V_s^2$ \\
4 & \code{lumped\_K}   & $K_s\,\tfrac12\rho V_s^2$, $K_s$ asserted \\
5 & \code{unstated}    & $0$, reported as unsized \\
\bottomrule
\end{tabular}
\end{center}
A line reports both its best and its \emph{weakest} method. In rank 3,
$L^{\mathrm{tube}}$ is the stated length with fitting bodies (less thread
engagement) removed when the length basis is ``overall''; if any fitting lacks
a body length the stated length is used whole and warned. A fitting with its
own $K$ uses it; one with its own bore is priced in that bore and referred back
to the segment's:
\begin{equation}
  K_j^{\mathrm{(seg)}} = K_j^{(D_j)}\Bigl(\frac{D_s}{D_j}\Bigr)^{4}.
\end{equation}
Bore changes between segments are \emph{derived} contractions/expansions
referred to the smaller bore.

\subsection{Orifices}
\paragraph{Fixed $C_d$ (\code{orifice/cd}).}
\begin{equation}
  \Delta p = \frac{\mdot^2}{2\rho\,(C_dA)^2},\qquad A=\tfrac{\pi}{4}D^2 .
  \label{eq:orifice}
\end{equation}
\paragraph{ISO 5167 (\code{orifice/iso5167}).} Downstream pressure from
\code{fluids.differential\_pressure\_meter\_solver} with Reader--Harris/Gallagher
$C(\beta,\mathrm{Re},\text{taps}=D)$ and expansibility $\epsilon\equiv1$
(incompressible). Infeasible flow raises rather than returning a negative
pressure.

\subsection{Compressible orifice (\code{gas\_orifice/isentropic})}\label{sec:gasorifice}
% NOTE: Omitted from revision 1. It is the model behind every relief orifice.
\begin{align}
  r^* &= \Bigl(\frac{2}{\gamma+1}\Bigr)^{\gamma/(\gamma-1)},\qquad
  \mdot^{\mathrm{ceil}} = C_dA\,p_1\sqrt{\frac{\gamma}{R_sT_1}}\Bigl(\frac{2}{\gamma+1}\Bigr)^{\frac{\gamma+1}{2(\gamma-1)}},\\
  \mdot(r) &= C_dA\,p_1\sqrt{\frac{2\gamma}{(\gamma-1)R_sT_1}\Bigl(r^{2/\gamma}-r^{(\gamma+1)/\gamma}\Bigr)},
  \qquad r=\frac{p_2}{p_1}\in(r^*,1).
\end{align}
The branch relation inverts $\mdot(r)$ for $r$ by Brent's method on
$[r^*,1)$, where it is monotone; $|\mdot|\ge\mdot^{\mathrm{ceil}}$ is
infeasible and the branch is written in choked form (\cref{eq:chokerow}) with
$\Delta p_{\mathrm{crit}} = p_1(1-r^*)$. Missing $\gamma$ or $R_s$ default to
$1.4$ and \SI{296.8}{\joule\per\kilogram\per\kelvin} (air/N$_2$).

\subsection{Manifold (\code{manifold/plenum})}\label{sec:manifold}
% NOTE: Omitted from revision 1. A COPV manifold with PT, relief, fill QD and
% regulator feed is the common case on this stand.
A plenum node plus one branch per port. A \emph{flow} port carries
\begin{equation}
  \Delta p_{\mathrm{port}} = \bigl(K_{\mathrm{turn}}+K_{\mathrm{fit}}\bigr)\tfrac12\rho V_{\mathrm{port}}^2,
  \qquad
  \Delta p^{\mathrm{static}} = -\rho g\,z_{\mathrm{port}},
\end{equation}
with $K_{\mathrm{turn}}=1.3$ by default (a branch tee at 90$^\circ$ is $\approx1.0$--$1.4$),
or derived as $K_{\mathrm{contraction}}(D_{\mathrm{plenum}}\!\to\!D_{\mathrm{port}}) + 1.3$
when the port is smaller than the plenum. \emph{Instrument} and \emph{conditional}
(shut relief) ports are declared dead ends and back-filled exactly.

\subsection{Valve (\code{valve})}\label{sec:valve}
\begin{equation}
  K_{\mathrm{valve}} = \frac{1.6\times10^{9}\,D^{4}}{\bigl(C_v^{\mathrm{eff}}/1.56\bigr)^{2}}\quad[D\ \text{in m}],
  \qquad
  C_v^{\mathrm{eff}}(\sigma) = \max\bigl(C_v\,\chi(\sigma),\ C_v^{\mathrm{leak}}\bigr),
  \qquad
  \Delta p = K_{\mathrm{valve}}\tfrac12\rho V^2,
  \label{eq:cvk}
\end{equation}
where $\chi$ is linear, equal-percentage, quick-opening (from
\code{fluids.control\_valve}) or a tabulated $C_v(\sigma)$ curve. Isolation
(\cref{sec:reduce}) is decided on travel $\sigma\le10^{-3}$, not on the
resulting $C_v$.

\subsection{Compressible flow through a valve, and choking}\label{sec:choking}
IEC 60534-2-1 sizes a valve on gas as
\begin{equation}
  W = N_6\,C\,Y\sqrt{x\,p_1\rho_1},\qquad
  x=\frac{p_1-p_2}{p_1},\qquad
  Y=1-\frac{x}{3F_\gamma x_T},\qquad
  F_\gamma=\frac{\gamma}{1.40},
  \label{eq:iec}
\end{equation}
with choking at $x_{\mathrm{crit}}=F_\gamma x_T$ where $Y=\tfrac23$ exactly.
The ceiling used in \cref{eq:chokerow} is
\begin{equation}
  \mdot^{\mathrm{ceil}} = \mathcal{C}\,C_v^{\mathrm{eff}}\,\frac{2}{3}
     \sqrt{F_\gamma x_T\,\frac{p_1}{1000}\,\rho_1},
  \qquad
  \text{choked}\iff \Delta p_{\mathrm{avail}}\ge F_\gamma x_T\,p_1,
  \label{eq:ceiling}
\end{equation}
with $x_T=0.7$ unless the valve declares one (its true range is
$\approx0.3$--$0.8$, and it is the largest single uncertainty in a choked-gas
prediction).

\begin{remark}[Calibration rather than transcription]
$\mathcal{C}$ depends on the unit set and on the $C_v$/$K_v$ pairing
($1\,K_v = 1.156\,C_v$). A transcribed $N_6$ was wrong by $10.5\times$, then by
exactly $1/0.865^2$. $\mathcal{C}$ is therefore obtained once at import by
inverting \code{fluids.control\_valve.size\_control\_valve\_g} at a reference
point; its invariance across He/N$_2$, \SIrange{200}{4500}{\psi},
\SIrange{250}{293}{\kelvin} and $C_v\in[1.6,60]$ to five figures is the check
that \cref{eq:ceiling} has the right \emph{form}.
\end{remark}

\paragraph{Liquid choking and cavitation.} $p_2^{\mathrm{choke}} = p_2(p_{\mathrm{sat}},p_{\mathrm{crit}},F_L,p_1)$
per IEC 60534; undefined (and so reported) when $p_{\mathrm{sat}}$ is unknown.
The cavitation index $\sigma_c=(p_1-p_v)/(p_1-p_2)$ is a diagnostic.

\subsection{Check valve (\code{check\_valve})}
\begin{equation}
  \Delta p = \begin{cases}
    K(C_v)\tfrac12\rho V^2 + \Delta p_{\mathrm{crack}}, & \mdot\ge0,\\
    K(C_v^{\mathrm{leak,rev}})\tfrac12\rho V^2, & \mdot<0,
  \end{cases}
\end{equation}
discontinuous at $\mdot=0$ by $\Delta p_{\mathrm{crack}}$, which is why a
stub behind an unopened check valve is reported indeterminate
(\cref{sec:reduce}).

\subsection{Measured element (\code{measured})}
$\Delta p = \mathcal{I}_{\mathrm{lin}}\bigl[\Delta p(\mdot)\bigr](|\mdot|)$,
linear interpolation of a flow-bench table. It refuses outside the measured
range rather than extrapolating; unlike the property chain there is no
defensible fallback.

\subsection{Pressure regulator (\code{regulator/droop})}\label{sec:regulator}
% EXPAND: a regulator defends a setpoint; it does not have a loss.
% NOTE: Say why there is no force balance (needs the diaphragm's effective
% area vs lift, which no vendor publishes).
The commanded setpoint is either fixed or dome-loaded,
\begin{equation}
  p_{\mathrm{set}} = \begin{cases} p_{\mathrm{dome}} + \Delta p_{\mathrm{bias}} & \text{dome-loaded (the 1092-50: } \Delta p_{\mathrm{bias}}=\SI{50}{\psi}),\\ p_{\mathrm{set}}^{\mathrm{hand}} & \text{otherwise,}\end{cases}
\end{equation}
where $p_{\mathrm{dome}}$ is a \emph{signal} (\cref{sec:pid}). The effective
outlet target is affine in supply pressure and flow:
\begin{equation}
  p^{\mathrm{eff}}_{\mathrm{set}}(\mdot,p_1)
   = p_{\mathrm{set}} + S\,(p_1^{\mathrm{ref}}-p_1) - D_r\frac{|\mdot|}{\mdot_{\mathrm{rated}}},
  \label{eq:reg}
\end{equation}
$S$ the supply-pressure coefficient in the datasheet's own units
(\SI{17}{\psi} per \SI{1000}{\psi} for the 1092-50; it carries a
\emph{pressure ratio} dimension so a bare $0.017$ is rejected), $D_r$ the droop
at rated flow. The branch relation is
\begin{equation}
  \Delta p_{\mathrm{reg}}(\mdot) = \max\Bigl(p_1-p^{\mathrm{eff}}_{\mathrm{set}},\ K_{\mathrm{seat}}\tfrac12\rho V^2,\ 0\Bigr),\qquad
  \Delta p_{\mathrm{reg}}(0)=\max\bigl(p_1-p_{\mathrm{lock}},0\bigr),\quad
  p_{\mathrm{lock}}=p_{\mathrm{set}}+\Delta p_{\mathrm{lock}},
\end{equation}
with $K_{\mathrm{seat}}$ from the seat $C_v$ via \cref{eq:cvk}. \emph{Saturated}
means the seat term binds. \emph{Shut} (choked form with ceiling $0$) means
downstream already exceeds the setpoint:
\begin{equation}
  \text{shut}\iff \Delta p_{\mathrm{avail}} < \bigl(p_1 - p^{\mathrm{eff}}_{\mathrm{set}}(0,p_1)\bigr) - 0.005\,p^{\mathrm{eff}}_{\mathrm{set}},
\end{equation}
the half-percent band existing so a tank sitting on its setpoint does not flip
the branch equation every step.

\begin{proposition}[Droop makes the branch determinate]
With $D_r=0$ the branch equation is satisfied for every $\mdot$ between two
fixed pressures; $J_{bb}\equiv0$ and the flow is undetermined. With $D_r>0$ the
row has slope $D_r/\mdot_{\mathrm{rated}}$ and $\mdot$ is unique. Zero droop
is therefore reported as a violation, and a perfect setpoint is a separate
named model (\code{ideal}).
\end{proposition}

\begin{remark}[Reverse flow]
\Cref{eq:signed} mirrored would place a step of $2(p_1-p_{\mathrm{set}})$ across
$\mdot=0$ and a finite-difference slope of order $10^{16}$. Reverse flow is
instead the poppet on its seat: $\Delta p^{\mathrm{tot}} = \max(p_1-p_{\mathrm{lock}},0) + R_{\mathrm{rev}}\,\mdot$
with $R_{\mathrm{rev}}=10^{10}\ \si{\pascal\second\per\kilogram}$ --- continuous
at zero, steep, and still descendable.
\end{remark}

Two further named models: \code{curve} reads outlet against flow from a
datasheet curve, shifted by the dome and corrected by $S$; \code{measured}
reads $\Delta p(\mdot)$ directly.

\subsection{Injector leg}\label{sec:injectorleg}
% NOTE: Omitted from revision 1 in its Re-dependent form. This must agree
% with EngineDesign at the same operating point or the two tools model
% different injectors.
\begin{equation}
  \Delta p = \frac{\mdot^2}{2\rho\,(C_d A)^2},\qquad
  C_d(\mathrm{Re}) = \mathrm{clip}\Bigl(C_d^{\infty}(d_h) - \frac{a_{\mathrm{Re}}}{\sqrt{\mathrm{Re}}},\ C_d^{\min},\ C_d^{\infty}(d_h)\Bigr),\qquad
  \mathrm{Re}=\frac{\rho V d_h}{\mu},\ V=\frac{|\mdot|}{\rho A},
  \label{eq:cd}
\end{equation}
with optional multiplicative pressure and temperature corrections
$1+a_p(p/p_{\mathrm{ref}}-1)$, $1+a_T(T/T_{\mathrm{ref}}-1)$, and a geometry
dependence
\begin{equation}
  C_d^{\infty}(d) = \mathrm{clip}\Bigl(
    \begin{cases} C_d^{\infty}\,(d/d_{\mathrm{ref}})^{0.20}, & d<d_{\mathrm{ref}}\\ C_d^{\infty} + 0.015\ln(d/d_{\mathrm{ref}}), & d\ge d_{\mathrm{ref}}\end{cases},\ 0.48,\ 0.62\Bigr),
  \quad d_{\mathrm{ref}}=\SI{2}{\milli\metre},
\end{equation}
when the engine config enables it. $A$ and $d_h$ are extracted per injector
type (impinging: $n$ round jets; pintle: discrete ox orifices, fuel annulus;
coaxial: core ports and outer annulus). Defaults when the config has no
discharge block: $C_d^\infty=0.6$, $C_d^{\min}=0.15$.

%==============================================================================
\section{Vessel models}\label{sec:vessels}
%==============================================================================

\subsection{State pairs}
% NOTE: The ullage is (m, U) because a real gas's energy balance is exact only
% in u; the liquid is (m, T) because liquid density is ~180x less sensitive to
% pressure, so recovering p from (rho, u) amplifies error by that factor.
\begin{equation}
  \text{gas: } (m,U,T_w),\ \ \rho=\frac{m}{V},\ u=\frac{U}{m},\ (p,T,h)=\mathrm{EOS}(\rho,u);
  \qquad
  \text{liquid: } (m_\ell, T_\ell),\ \ \rho_\ell = \rho_{\mathrm{sat}}(T_\ell).
\end{equation}
A $(\rho,u)$ inversion that fails is raised as \emph{ullage condensed}: either
a long unpressurised hold has genuinely chilled the gas into a liquid-like
state, or the step is too large.

\subsection{Fixed-volume gas vessel (bottle, plenum)}
\begin{align}
  \dot m &= \mdot_{\mathrm{in}}-\mdot_{\mathrm{out}},\\
  \dot U &= \mdot_{\mathrm{in}}h_{\mathrm{in}} - \mdot_{\mathrm{out}}\,h(\rho,u) + \Qdot_{\mathrm{ext}} + \Qdot_w,
  \qquad \Qdot_w = (hA)_w\,(T_w-T_g),\qquad
  \dot T_w = -\frac{\Qdot_w}{m_wc_w}.
  \label{eq:bottle}
\end{align}
Outflow carries the vessel's own \emph{enthalpy}, not internal energy; the
difference is the flow work that cools what remains.

The wall term carries a stirring factor $\sigma_s\ge 1$ while gas is
being charged in, $\Qdot_w = \sigma_s\,hA\,(T_w-T)$: $hA$ is the still-gas
conductance and a charge jet drives forced convection several times it. With
$\sigma_s=1$ a \SI{25}{\second} GSE fill of \SI{10}{\kilogram} of N$_2$ into
\SI{44}{\litre} is an adiabatic charge and lands at \SI{380}{\kelvin}
($\gamma T_{\mathrm{bank}}$ less what the wall takes), then sags \SI{12}{\percent}
in a minute; the stand does not, so the default is $\sigma_s=20$
(\SI{336}{\kelvin} over a \SI{316}{\kelvin} wall, \SI{3}{\percent}), an
order-of-magnitude estimate to be calibrated against the bottle RTD. The same
factor applies to an ullage while it is being pressed --- gas arriving with
the tank more than \SI{2}{\percent} below its supply --- and not to the few
grams a second two tanks at lockup trade through a shared manifold, which is
slosh, not a jet (applied to that, it ratcheted both tanks \SI{16}{\psi}
above lockup in ten seconds).
% NOTE: Replace 20 with the value the RTD gives after a real GN2 High Press.

\subsection{Propellant tank}
\begin{align}
  \dot V_u &= \frac{\mdot_\ell}{\rho_\ell},\qquad \Pi = p\,\dot V_u,\\
  \dot U_u &= \mdot_g h_g + \Qdot_{\mathrm{ext}} + \Qdot_w - \Qdot_i - \Pi, \label{eq:ullage}\\
  \Qdot_{\mathrm{wet}} &= U_{\mathrm{wet}}\,A_{\mathrm{wet}}(\ell)\,(T_{w,\ell} - T_\ell)\ \ (\text{chilldown; 0 unless requested}),\label{eq:chilldown}\\
  \dot T_w &= -\frac{\Qdot_w + \Qdot_{\mathrm{wet}}}{m_wc_w}\ \ (\text{one wall node}),\qquad
  \dot T_{w,u} = -\frac{\Qdot_w}{(1-f_{\mathrm{wet}})\,m_wc_w},\ \ 
  \dot T_{w,\ell} = -\frac{\Qdot_{\mathrm{wet}}}{f_{\mathrm{wet}}\,m_wc_w}\ \ (\text{two}),\label{eq:twowall}\\
  \dot m_\ell &= -\mdot_\ell - \mdot_v - \mdot_{v,w},\qquad
  \dot T_\ell = \frac{\Qdot_i + \Qdot_{\mathrm{wet}} - \Qdot_{\mathrm{boil}} - \Qdot^{\mathrm{lat}}}{m_\ell c_{p,\ell}},\qquad
  \dot t_c = 1,\\
  \Qdot_{\mathrm{amb},u} &= U_{\mathrm{amb}}\,(A_H - A_{\mathrm{wet}}(\ell))\,(T_{\mathrm{amb}}-T_{w,u}),\qquad
  \Qdot_{\mathrm{amb},\ell} = U_{\mathrm{amb}}\,A_{\mathrm{wet}}(\ell)\,(T_{\mathrm{amb}}-T_{w,\ell}),\label{eq:leak}\\
  \mdot_{v,w} &= \frac{\Qdot_{\mathrm{boil}}}{h_{fg}(T_\ell)},\qquad
  \Qdot_{\mathrm{boil}} = \begin{cases}\Qdot_{\mathrm{wet}} & T_{w,\ell} > T_{\mathrm{sat}}(p) + \Delta T_{\mathrm{onb}}\\ 0 & \text{otherwise},\end{cases}\label{eq:wallboil}
\end{align}
with the wetted-wall conductance itself regime-dependent,
\begin{equation}
  h_{\mathrm{wet}} = \begin{cases} h_{\mathrm{film}} & T_{w,\ell} - T_{\mathrm{sat}}(p) \ge \Delta T_{\mathrm{Leid}} \\ h_{\mathrm{nuc}} & 0 < T_{w,\ell} - T_{\mathrm{sat}}(p) < \Delta T_{\mathrm{Leid}} \\ h_{\mathrm{film}} & \text{otherwise (single phase)},\end{cases}
  \label{eq:regimes}
\end{equation}
$h_{\mathrm{film}}=\SI{100}{\watt\per\square\metre\per\kelvin}$ (a vapour film
insulates a wall far above saturation), $h_{\mathrm{nuc}}=\SI{3000}{\watt\per\square\metre\per\kelvin}$
once the superheat is under the Leidenfrost point $\Delta T_{\mathrm{Leid}}=\SI{40}{\kelvin}$
and the liquid wets the metal, and $\Delta T_{\mathrm{onb}}=\SI{2}{\kelvin}$ the incipience
superheat below which the wall warms the liquid rather than boiling it. All three
are estimates from cryogen boiling curves and are dialled from the cockpit's
Configuration tab.

\paragraph{Stratification.} A quiescent cryogen does not warm as a lump. Heat
that reaches the wetted wall and does not boil rides up the wall in a
natural-convection boundary layer and pools at the surface; the bulk below stays
cold, and the ullage sees the surface. The tank therefore carries a surface
temperature $T_s$ apart from the bulk $T_\ell$: a slab of depth $\delta_s$
over the interface, $m_s = \rho_\ell A_i \delta_s$ (never more than the liquid
present), with
\begin{equation}
  m_s c_\ell \dot T_s = \Qdot_i + \Qdot_{\mathrm{wet}} - \Qdot_{\mathrm{boil}} - \Qdot^{\mathrm{lat}} - h_{\mathrm{mix}} A_i (T_s - T_\ell),
  \qquad
  (m_\ell - m_s) c_\ell \dot T_\ell = h_{\mathrm{mix}} A_i (T_s - T_\ell),
  \label{eq:strat}
\end{equation}
and the collapse and vapour closures (\cref{sec:collapse,sec:vapour}) evaluated at
$T_s$. The wall's convective share reaches the layer only in the proportion
$[(T_{w,\ell}-T_s)/(T_{w,\ell}-T_\ell)]_0^1$: liquid warmed at a wall arrives no hotter
than the wall, so a layer at the wall's temperature takes nothing more from it and
the remainder goes to the bulk (in a slab of fixed depth, ``thickening the layer''
is warming the bulk). The ullage's heat needs no such bound; the collapse term
already follows $T_g - T_s$. The layer itself is bounded at saturation for the
tank's total pressure (and two kelvin under the critical temperature, where the
saturation line ends): a surface there boils rather than superheats, and heat
that would have warmed it evaporates it instead. Without that a room-temperature
dry wall over a small ullage ran the surface to \SI{154.6}{\kelvin} and every
property call failed. What sustains a climb of tens of psi a minute is
therefore warm hardware -- the dry wall the liquid never touched, still near
\SI{200}{\kelvin} after a two-minute load -- warming the ullage, which warms
the surface. $\delta_s=\SI{1}{\centi\metre}$ and $h_{\mathrm{mix}}=\SI{5}{\watt\per\square\metre\per\kelvin}$
by default (estimates; the tank's own shut-vent trace calibrates them); a load
resets $T_s$ to $T_\ell$ since a pour stirs the liquid. Off ($\delta_s=0$, the
library default) the liquid is well mixed and \cref{eq:strat} collapses to the
single bulk balance. The three closures together are what set a shut LOX tank's
climb: boiling the whole leak into the ullage gives $\sim\SI{100}{psi/min}$,
warming fifteen kilograms gives $\sim\SI{0.3}{psi/min}$, and a centimetre of
warm surface under a \SI{75}{\watt} leak gives the tens of psi per minute a stand
actually shows.
with \cref{eq:leak} added to the respective wall lumps of \cref{eq:twowall}
($U_{\mathrm{amb}}=0$ by default in the library; on the cockpit
$U_{\mathrm{amb}} = (1/h_{\mathrm{air}} + t_{\mathrm{ins}}/k_{\mathrm{ins}})^{-1}$ with
$h_{\mathrm{air}}=\SI{8}{\watt\per\square\metre\per\kelvin}$ and the insulation the
drawing declares --- an inch of fiberglass on the LOX tank gives
\SI{1.3}{\watt\per\square\metre\per\kelvin}, about \SI{75}{\watt}; a bare tank sees
$h_{\mathrm{air}}$, six times that) and \cref{eq:wallboil}
the wall-boiling closure (off by default in the library; on with the vapour
model on the cockpit). Nucleation needs the wall above saturation at the
tank's \emph{total} pressure, so a LOX tank at atmosphere with its vent shut
boils and climbs --- self-limiting, as the rising pressure lifts
$T_{\mathrm{sat}}$ toward the wall --- while the same tank under
\SI{38}{\bar} of pressurant is a hundred kelvin subcooled and the leak warms
the bulk instead. Gas leaving through a vent takes pressurant and vapour in
proportion to their masses --- only along paths that end at a vent boundary,
found by following the solved flows through the graph and splitting at each
junction; what goes to a neighbouring tank through a shared manifold and back
is booked as pressurant both ways, since the network carries no species. The
vapour term was missing until 2026-09-11 and a venting LOX tank kept every
gram it boiled; applied to the manifold slosh as well, it ratcheted a pair of
tanks at lockup upward.
with $\Qdot_i$ from \cref{sec:collapse}, $\mdot_v$ and $\Qdot^{\mathrm{lat}}$
from \cref{sec:vapour}, and $t_c$ the interface contact time, a state so a
repressurisation can reset it. The wall may be one lump (the default, and
every state built before the second node existed) or two: the wall under the
liquid, $T_{w,\ell}$, and the wall above it, $T_{w,u}$, split by the wetted-area
fraction $f_{\mathrm{wet}} = A_{\mathrm{wet}}(\ell)/A_{\mathrm{wet}}(H)$ of the
tank as it stands, with no conduction between them (a thin wall over a tank's
height is a twenty-minute time constant).
\begin{remark}[Why two nodes]
One lump cannot describe a loaded cryogenic tank. At the liquid's temperature
it condenses the pressurant on the ullage face --- N$_2$ at \SI{39}{\bar}
saturates at \SI{113}{\kelvin}, well above a \SI{90}{\kelvin} wall --- and at
ambient it boils the liquid at kilowatts through $U_{\mathrm{wet}}$. Primed as
freshly filled (wall \SI{293}{\kelvin}, $t_c=\SI{2}{\second}$), a LOX tank with
collapse and chilldown enabled sat \SI{139}{\psi} under lockup at $T$-0 and its
recovery read as an ignition transient. A primed tank now carries
$T_{w,\ell}=T_\ell$, $T_{w,u}=T_g$, and $t_c=$ the stated pad hold.
\end{remark}
\begin{remark}[Initial vapour]
A tank primed with a vapour model fitted, liquid present and $t_c>0$ starts
with its propellant's equilibrium vapour, $p_{\mathrm{vap}}=p_{\mathrm{sat}}(T_\ell)$
(only where $p_{\mathrm{sat}}<\tfrac12 p$), and the pressurant makes up the
\emph{remainder} of the stated total. Built dry, a saturated cryogen boils
that vapour off after the regulator has locked up at the total, and nothing
can take the extra \SI{15}{\psi} (LOX at \SI{90}{\kelvin}) back out; the trace
then opens above lockup. Default and \code{none}: no vapour, exactly as before.
\end{remark} Pressures seen by the network:
\begin{equation}
  p_{\mathrm{ullage}} = p_{\mathrm{EOS}}(\rho_g,u_g) + p_{\mathrm{vap}},\qquad
  p_{\mathrm{outlet}} = p_{\mathrm{ullage}} + \rho_\ell\,g\,\ell(V_\ell),
  \label{eq:outlet}
\end{equation}
where the level $\ell$ inverts $V_{\mathrm{below}}(\ell)$ by Brent's method.

\paragraph{Simultaneous in- and out-flow.} The cockpit feeds \cref{eq:ullage}
one net gas stream with an effective enthalpy,
\begin{equation}
  \mdot_g = \mdot_{g,\mathrm{in}}-\mdot_{g,\mathrm{out}},\qquad
  h_g = \frac{\mdot_{g,\mathrm{in}}h_{\mathrm{in}} - \mdot_{g,\mathrm{out}}h(\rho_g,u_g)}{\mdot_g},
\end{equation}
so a tank being pressed and vented at once does not price the whole exchange
at the supply's enthalpy.

\subsection{Geometry}
Cylinder radius $r$, barrel $L_b$, semi-ellipsoidal heads of depth $c=r/\text{ratio}$
(2:1 default; hemispherical $c=r$; flat $c=0$):
\begin{gather}
  V_{\mathrm{head}}(d) = \frac{\pi r^2 d^2(3c-d)}{3c^2},\qquad
  V_{\mathrm{head}}^{\mathrm{tot}}=\tfrac23\pi r^2c,\qquad
  A_{\mathrm{xs}}(d) = \pi r^2\Bigl[1-\Bigl(\frac{c-d}{c}\Bigr)^2\Bigr],\\
  A_{\mathrm{wet}}^{\mathrm{head}}(d) = 2\pi r d\ \ (\text{hat-box; exact for a hemisphere, a few percent low for 2:1}),\qquad
  A_{\mathrm{wet}}^{\mathrm{barrel}}=2\pi r\,\ell_{\mathrm{barrel}}.
\end{gather}
$A_i = A_{\mathrm{xs}}(\ell)$ is the interface area of \cref{sec:collapse}.
A tabulated $V(h)$ (with optional $A(h)$, $A_{\mathrm{wet}}(h)$) from CAD is the
escape hatch. The cockpit builds a tank from the drawing's volume and diameter,
2:1 heads, solving $L_b = \max\bigl((V-\tfrac43\pi r^2\cdot\tfrac r2)/(\pi r^2),0\bigr)$.

\subsection{Ullage collapse: transient interfacial conduction}\label{sec:collapse}
% NOTE: Why not a steady h: the layer depth grows as sqrt(pi alpha t) -- 1.1 mm
% at 5 s, 3.8 mm at 60 s for LOX -- so no single coefficient fits both a burn
% and a pad hold.
\begin{equation}
  \Qdot_i = \eta_c\,A_i\,(T_g-T_\ell)\,\frac{\beta}{\sqrt{\pi\,\max(t_c,\,t_{\min})}},
  \qquad \beta=\sqrt{k\rho c_p}\Big|_{\mathrm{sat},T_\ell},\qquad t_{\min}=\SI{1}{\milli\second},
  \label{eq:collapse}
\end{equation}
whose closed-form integral at constant $A_i,\Delta T$ is $Q_i(t)=2\eta_cA_i\Delta T\beta\sqrt{t/\pi}$.
$\eta_c=1$ is quiescent conduction; a fitted enhancement from a cold flow goes
there and nowhere else. Named alternatives: \code{none}.
\begin{remark}
Quiescent conduction is a \textbf{lower bound}: free convection, slosh, jet
scrubbing and N$_2$-over-LOX condensation all move more. A pressurant budget on
\cref{eq:collapse} is optimistic.
\end{remark}

\subsection{Propellant vapour}\label{sec:vapour}
Dalton partial pressures, \emph{not} a mixture EOS:
\begin{equation}
  p_{\mathrm{vap}} = \begin{cases}
    p^{\ell}_{\mathrm{EOS}}\bigl(\rho=m_v/V_u,\ T=T_g\bigr) & \text{if }>0,\\
    \dfrac{m_v R_s^{\ell}T_g}{V_u} & \text{otherwise (ideal-gas fallback; } R_s^\ell = R/M_\ell),
  \end{cases}
\end{equation}
the fallback existing because the Helmholtz formulation returns non-physical
values at very low partial density (\SI{1}{\micro\gram} of O$_2$ in
\SI{0.5}{\litre} gave \SI{-6057}{\psi}). Evaporation (\code{saturated}):
\begin{equation}
  \psi = \mathrm{clip}\Bigl(\frac{p_{\mathrm{sat}}(T_\ell)-p_{\mathrm{vap}}}{p_{\mathrm{sat}}(T_\ell)},\,-1,\,1\Bigr),\qquad
  \Qdot^{\mathrm{lat}} = f_{\mathrm{sat}}\,\psi\,\bigl(\Qdot_i+\Qdot_{\mathrm{wet}}\bigr),\qquad
  \mdot_v = \frac{\Qdot^{\mathrm{lat}}}{h_{fg}(T_\ell)},
  \label{eq:evap}
\end{equation}
$h_{fg}=h_{\mathrm{sat}}^v-h_{\mathrm{sat}}^\ell$ from the same EOS, $f_{\mathrm{sat}}=1$
by default. $\psi<0$ is condensation. Default model: \code{none} in the library; the cockpit fits \code{saturated} with wall boiling (\cref{eq:wallboil}) and the ambient leak (\cref{eq:leak}) on, since 2026-09-11.

\begin{proposition}[Vapour carries mass and partial pressure, not energy]\label{prop:reference}
$u_g$ is referred to the pressurant's reference state and $h_{\mathrm{sat}}$
to the propellant's; these are unrelated (saturated LOX at \SI{90}{\kelvin}:
\SI{-133.7}{\kilo\joule\per\kilogram}; helium at ullage conditions:
\SI{+628.6}{\kilo\joule\per\kilogram}). Adding $\mdot_vh_{\mathrm{sat}}$ into
\cref{eq:ullage} is a category error that drives the state negative within a
step. The energy that produced the vapour has already left \cref{eq:ullage}
through $\Qdot_i$; the vapour contributes no energy term. The stated cost is
that the vapour's own sensible energy is not returned to the ullage, biasing it
cool.
\end{proposition}

\subsection{Pressurant budget (derived quantity)}
For a stated expulsion the gas inflow that holds $p_{\mathrm{ullage}}$ constant
is found per step by bisection on $\mdot_g$ (monotone), integrating
\cref{eq:ullage} with expansion, interfacial and wall terms; the result is
reported with the residual pressure error rather than trusted.

%==============================================================================
\section{Thermal transport in lines}\label{sec:thermal}
%==============================================================================

\subsection{Enthalpy propagation and Joule--Thomson}
% NOTE: Emphasise that temperature is *not* a solver unknown. The walk runs
% after each solve and lags it by one step -- deliberately, so it cannot
% perturb Newton.
An adiabatic component with no shaft work conserves specific enthalpy. After
each network solve, node temperatures are updated by a fixed-point walk:
vessel nodes are seeded from their integrated states; then, for up to 12 sweeps
or until no node moves by more than \SI{0.05}{\kelvin},
\begin{equation}
  h_i = \frac{\sum_{b\in\mathcal{B}^{\mathrm{in}}_i}|\mdot_b|\,h_b^{\mathrm{out}}}{\sum_{b\in\mathcal{B}^{\mathrm{in}}_i}|\mdot_b|},
  \qquad
  T_i = T_{\mathrm{EOS}}(p_i,h_i),
  \qquad
  \mathcal{B}^{\mathrm{in}}_i = \{b: \text{flow enters } i,\ |\mdot_b|>\SI{e-6}{\kilogram\per\second}\},
  \label{eq:enthalpy}
\end{equation}
with $h_b^{\mathrm{out}} = h(p_{\mathrm{up}},T_{\mathrm{up}}) + \Qdot_{w,b}/|\mdot_b|$
(\cref{sec:wall}). Direction follows the sign of $\mdot_b$. Vessel nodes are
never overwritten, but the enthalpy arriving at an ullage node, $h_i$ of
\cref{eq:enthalpy}, is recorded and is $h_{\mathrm{in}}$ in \cref{eq:ullage}
whenever anything arrived; only a node nothing reached that step falls back to
the bottle's enthalpy.
\begin{remark}[The source matters, not just the path]
An earlier revision read $h_i$ only with line walls enabled and quoted the
bottle otherwise, on the argument that an adiabatic path conserves enthalpy.
That holds for gas that \emph{came from the bottle}. Two tanks primed on one
press manifold at equal pressure trade $\mathcal{O}(\si{\gram\per\second})$
back and forth each step; gas arriving from the other tank's \SI{293}{\kelvin}
ullage was priced at $h(\SI{4500}{\psi},\SI{293}{\kelvin})$, which for helium
is \emph{above} $h(\SI{550}{\psi},\SI{293}{\kelvin})$, so every exchange pumped
energy in: both ullages warmed \SI{6}{\kelvin} and rose \SI{15}{\psi} above
lockup with no flow through the regulator. Nitrogen ran the same error the
other way and cooled. Reading $h_i$ unconditionally changes nothing for
bottle-fed gas beyond the walk's \SI{0.05}{\kelvin} tolerance.
\end{remark}
\begin{remark}
\Cref{eq:enthalpy} produces Joule--Thomson at every throttle with no
gas-specific code: N$_2$ throttled \SIrange{4500}{550}{\psi} cools
${\approx}\,\SI{30}{\kelvin}$; He (inversion ${\approx}\,\SI{45}{\kelvin}$) warms
${\approx}\,\SI{15}{\kelvin}$. A fixed delivery temperature erases the difference.
\end{remark}

\subsection{Line walls: NTU--effectiveness exchange}\label{sec:wall}
% EXPAND: the icicle observation, and that it is evidence for the metal's own
% thermal mass (path A) and against the room (path B, 0.04-0.18 K).
\begin{align}
  \mathrm{Nu} &= \begin{cases}3.66, & \mathrm{Re}<\mathrm{Re}_{\mathrm{lam}},\\ 0.023\,\mathrm{Re}^{0.8}\mathrm{Pr}^{\,n}, & \mathrm{Re}\ge\mathrm{Re}_{\mathrm{lam}},\ \ n=0.4\ (T_w>T_{\mathrm{in}}),\ 0.3\ (T_w<T_{\mathrm{in}}),\end{cases}\\
  \bar h &= \frac{\mathrm{Nu}\,k}{D},\qquad
  \mathrm{NTU}=\frac{\bar h A_s}{|\mdot|c_p},\qquad
  \eff = 1-e^{-\min(\mathrm{NTU},40)},\qquad
  A_s = \pi D\max(L,D),\label{eq:ntu}\\
  T_{\mathrm{out}} &= T_{\mathrm{in}}+\eff\,(T_w-T_{\mathrm{in}}),\qquad
  \Qdot_w = |\mdot|c_p\,\eff\,(T_w-T_{\mathrm{in}}).
\end{align}
Properties are evaluated at $(p_{\mathrm{up}},T_{\mathrm{in}})$. Zero flow gives
exactly zero exchange. $\mathrm{Nu}$ is checked against
\code{ht.turbulent\_Dittus\_Boelter} (agreement \SI{0.00}{\percent}) and
\code{ht.Nu\_conv\_internal} (\SI{1.1}{\percent}) in the benchmark.

\paragraph{Wall thermal state.} Once per solve (not per sweep),
\begin{equation}
  T_w^{+} = T_w - \frac{\Qdot_w\,\Delta t}{m_w\,c_w(T_w)},\qquad
  \text{then } T_w^{+}\leftarrow T_{\mathrm{in}}\ \text{if } (T_w-T_{\mathrm{in}})(T_w^{+}-T_{\mathrm{in}})<0,
\end{equation}
the clamp forbidding the wall from crossing the stream in one step (sign-agnostic:
warm gas through a cold LOX line is the same equation reversed). Metal:
\begin{equation}
  m_w = \tfrac{\pi}{4}\bigl[(D+2\tau)^2-D^2\bigr]L\rho_s + m_{\mathrm{fit}},\qquad
  m_{\mathrm{fit}} = \begin{cases} m_{\mathrm{fit}}^{\mathrm{weighed}} & \text{if given},\\ N\cdot\SI{113}{\gram}\,(D/\SI{10.92}{\milli\metre})^2 & \text{from a count},\end{cases}
\end{equation}
$\rho_s=\SI{8000}{\kilogram\per\cubic\metre}$; $c_w(T)$ piecewise-linear through
NIST values for austenitic stainless, $c_w(293)/c_w(77)\approx2.6$. The
bore-scaling is fitted over \SIrange{0.25}{0.5}{\inch} unions (measured
$32$, $62$, \SI{113}{\gram}; exponent $1.82$).

\subsection{Wall initial condition}
\begin{assumption}[No soak]\label{as:soak}
No heat transfer at zero flow; no model of how a stand drifts toward ambient
while it sits.
\end{assumption}
Each wall is seeded on the first solve for which \cref{eq:enthalpy} has
converged, to the temperature of its upstream node at that instant:
liquid lines at liquid temperature, lines off an ullage at the ullage's,
pressurant lines at the bottle's. Seeding from declared drawing values would put
a vent line off the top of a LOX tank \SI{200}{\kelvin} too cold.

\subsection{Validity of the lumped wall}\label{sec:biot}
$\mathrm{Bi}=\bar hL_c/k_s$ with $\bar h\approx\SI{2800}{\watt\per\square\metre\per\kelvin}$,
$k_s\approx\SI{15}{\watt\per\metre\per\kelvin}$:
\begin{center}
\begin{tabular}{lccl}
\toprule
$L_c$ & $\mathrm{Bi}$ & Over-credit at \SI{6}{\second} (vs 1-D conduction) & \\
\midrule
\SI{0.889}{\milli\metre} tube wall & 0.17 & $0\%$ & lumped exact \\
\SI{5}{\milli\metre} fitting body & 0.93 & $19\%$ & invalid \\
\SI{8}{\milli\metre} fitting body & 1.49 & $32\%$ & invalid \\
\bottomrule
\end{tabular}
\end{center}
Steel's diffusion depth is \SI{4.7}{\milli\metre} at \SI{6}{\second}. Carried
knowingly because the result is film-limited (\cref{sec:limits}).

%==============================================================================
\section{Engine coupling}\label{sec:engine}
%==============================================================================

\subsection{Chamber}
\begin{equation}
  p_c = \max\Bigl(\frac{(\mdot_{\mathrm{ox}}+\mdot_f)\,\eta_{c^*}\,c^*(p_c,\mathrm{MR})}{A_t},\ p_{\mathrm{amb}}\Bigr),
  \qquad \mathrm{MR}=\frac{\mdot_{\mathrm{ox}}}{\mdot_f},
  \label{eq:pc}
\end{equation}
iterated three fixed-point passes (the $p_c$-dependence of $c^*$ is weak). A
cold chamber sits at $p_{\mathrm{amb}}$ through its own nozzle; this is physics,
not a guard. Thrust and impulse, reported only when $p_c>1.001\,p_{\mathrm{amb}}$:
\begin{equation}
  F = \eta_n\,C_F(p_c,\mathrm{MR})\,p_c\,A_t,\qquad
  I_{sp}=\frac{F}{(\mdot_{\mathrm{ox}}+\mdot_f)\,g_0},\qquad
  \tau_{\mathrm{fill}} = \frac{V_c}{A_tc^*}\ \ (\text{reported, not integrated}).
\end{equation}
$\eta_{c^*}$ and $\eta_n$ are $1$ unless the engine config states them.

\subsection{Combustion closure}
$c^*,C_F,T_c,\gamma$ by trilinear interpolation (\code{RegularGridInterpolator})
of EngineDesign's cached CEA table on $(p_c,\mathrm{MR},A_e/A_t)$ at the
table's design expansion ratio; queries outside the $(p_c,\mathrm{MR})$ range
are \emph{clamped} to the edge and flagged, so a startup can run and a plateau
that is a table edge is identifiable.

\subsection{Closing the loop around the network}\label{sec:engineloop}
The chamber node is a fixed-pressure boundary whose value depends on the flows
it receives. Let $g(p)$ be the chamber pressure implied by solving the network
with the chamber pinned at $p$. The fixed point $f(p)=g(p)-p=0$ is found:
\begin{itemize}[nosep]
\item \textbf{Batch (transient, run):} by secant on $f$, with the secant step
      taken only when the local slope $f'<0.5$ (the map is nearly affine and
      contractive), otherwise a relaxed step $\omega f$ ($\omega=0.5$
      transient, $0.4$ run); floored at $p_{\mathrm{amb}}$; tolerance $10^{-5}$
      ($10^{-4}$ run) relative; 40 (30) iterations; non-convergence is
      \emph{reported}, since it is a design finding (an injector too soft for
      its feed system).
\item \textbf{Cockpit:} solved, not relaxed, on every coupling step. The
      root of $f$ is bracketed by construction --- $f(p_{\mathrm{amb}})\ge 0$
      whenever propellant arrives, $f(p_{\mathrm{feed}})<0$ where nothing
      flows --- and regula falsi with a bisection guard finds it to
      \SI{0.5}{\psi} in a handful of network solves, warm-started so a step
      on which the previous answer still holds costs no extra solve.
\end{itemize}
\begin{remark}[Why relaxation was abandoned]
The cockpit used $p_c\leftarrow p_c+\omega\,(g(p_c)-p_c)$ with $\omega=0.3$.
That is a fixed-point iteration on a map whose slope at the operating point
is $g'(p_c) = -p_c/(2\,\Delta p_{\mathrm{inj}})$, so the relaxed iteration
has multiplier $1-\omega+\omega g'$. For the 7000 N doublet
($\Delta p_{\mathrm{inj}}\approx\SI{66}{\psi}$ at $p_c\approx\SI{400}{\psia}$,
$g'\approx-3$) that is $-0.2$: stable. For the 7200 N doublet at
$\Delta p_{\mathrm{inj}}\approx\SI{40}{\psi}$ it is below $-1$, and the burn
flip-flopped frame by frame between \SI{718}{\psia} --- above the tank
pressure --- and nothing. The dynamics also depended on the tick length,
because the factor was per step. Solving the fixed point removes both.
\end{remark}
Flow into the chamber is clamped non-negative (a chamber does not receive
backflow; a negative mid-iteration means the injector cannot pass at that
$p_c$). Without an engine model the chamber boundary is the drawing's design
pressure while propellant is arriving (${>}\,\SI{1}{\gram\per\second}$) and
ambient otherwise.

\subsection{Mixture-ratio decomposition}
% NOTE: This is an identity, not a fit -- say so, and say what each factor is
% *for* (redrill vs re-plumb).
From \cref{eq:orifice} on each side,
\begin{equation}
  \mathrm{MR} = \underbrace{\frac{C_{d,\mathrm{ox}}A_{\mathrm{ox}}\sqrt{\rho_{\mathrm{ox}}}}{C_{d,f}A_f\sqrt{\rho_f}}}_{\text{face ratio}}
  \cdot
  \underbrace{\sqrt{\frac{\Delta p_{\mathrm{inj,ox}}}{\Delta p_{\mathrm{inj},f}}}}_{\text{feed term}},
\end{equation}
exactly, with $C_d$ \emph{recovered} from the solved $(\mdot,\rho,A,\Delta p)$
so the identity closes on the solver's own numbers and its residual measures
inconsistency rather than the difference between two $C_d$ models. Stiffness
$\Lambda_s=\Delta p_{\mathrm{inj},s}/p_c$ is reported per side against the
engine's own design band, else $0.20$ (soft) and $0.10$ (barely metering). The
fuel-side trim pressure to restore the design MR at the present flow is
$\Delta p_{\mathrm{inj,ox}}/(\mathrm{MR}_{\mathrm{des}}/\text{face})^2 - \Delta p_{\mathrm{inj},f}$,
an under-estimate because the fuel line's own loss rises with the flow.

%==============================================================================
\section{Transient formulation}\label{sec:dae}
%==============================================================================

\subsection{Semi-explicit index-1 DAE}
\begin{equation}
  \dot{\mathbf{x}} = \mathbf{f}(\mathbf{x},\mathbf{y},t),\qquad
  \mathbf{0}=\mathbf{g}(\mathbf{x},\mathbf{y},t),
  \label{eq:dae}
\end{equation}
$\mathbf{x}$ the vessel states (bottle: $m,U,T_w$; tank: $m_u,U_u,T_w,m_\ell,T_\ell,t_c$),
$\mathbf{g}$ from \cref{eq:mass,eq:branch}. Index 1 because
$\partial\mathbf{g}/\partial\mathbf{y}$ is nonsingular under \cref{sec:regular}.
Vessel states set the fixed boundary pressures (\cref{eq:outlet}); solved flows
return as vessel in/outflows, signed by topology. Gas arriving at a tank
carries the enthalpy of the \emph{upstream} node's $(p,T)$.

\subsection{Nested solution}
% EXPAND: the trade against monolithic (carrying p as algebraic unknowns).
At every right-hand-side evaluation $\mathbf{g}=\mathbf{0}$ is solved to
convergence. The constraint is satisfied exactly rather than to a tolerance,
the integrator state stays small, and a failure is localised to a named time
with a labelled state vector.

\subsection{Quasi-steady lines}\label{sec:qs}
\begin{assumption}[Quasi-steady]\label{as:qs}
Flow in every branch establishes much faster than the vessels change.
\end{assumption}
With mass-flow inertance $I=L/A$ (\si{\pascal\second\squared\per\kilogram}) and the local
quadratic-loss slope $\partial\Delta p/\partial\mdot\approx2\Delta p/|\mdot|$,
\begin{equation}
  \tau_{\mathrm{inert}} = \frac{I}{\partial\Delta p/\partial\mdot}
    = \frac{L\,|\mdot|}{2A\,\Delta p},
  \label{eq:inert}
\end{equation}
reported per branch \textbf{at that branch's own peak flow} (it goes as
$1/\mdot$, so a value read while a valve closes is large and meaningless).
\SI{2}{\metre} of \SI{0.375}{\inch} LOX line at \SI{1.5}{\kilogram\per\second}:
${\approx}\,\SI{9}{\milli\second}$. A $\tau_{\mathrm{inert}}$ approaching the
transient's timescale means the formulation does not apply to that branch; no
step change fixes it.

\subsection{Batch integration}
Stiff by construction (a \SI{40}{\milli\second} valve beside a \SI{5}{\second}
blowdown): BDF (default), Radau or LSODA; $\mathrm{rtol}=10^{-6}$,
$\mathrm{atol}=10^{-9}$; inner network tolerance $10^{-5}$. Command edges (each
command's start and end) are step boundaries, and the maximum step is bounded
by a quarter of the shortest gap between edges. Commanded actuators move by a
named shape --- smoothstep $3\xi^2-2\xi^3$ (default), linear, or step --- over
a declared travel time; a later command on the same signal supersedes an
earlier one from wherever the actuator has got to.

\subsection{Conservation audit}
\begin{equation}
  \mathcal{R} = \Bigl(\int\!\mdot_{\mathrm{in}}\,\mathrm{d}t - \int\!\mdot_{\mathrm{out}}\,\mathrm{d}t\Bigr) - \Delta\Bigl(\sum_{\text{state slots named }\ast\text{mass}}\Bigr),
\end{equation}
boundary flows by the trapezoid rule over the \emph{recorded} samples on
branches with exactly one end at a vessel node, inventory read from the state
vector directly. Two error sources: integrator tolerance (smooth blowdown) or
output quadrature (fast valve inside a slow burn); reported as a number so the
binding one can be identified.

\subsection{Explicit cockpit coupling}\label{sec:rc}
% NOTE: Say plainly that this is a different scheme from the batch path --
% forward Euler on the vessels with the network re-solved between sub-steps --
% and that its stability bound is the thing this subsection derives.
The interactive path advances vessels by forward Euler with the network
re-solved between sub-steps. A tick $\Delta t\le\SI{0.25}{\second}$ is divided
into $n_c$ coupling steps,
\begin{equation}
  n_c = \min\Bigl(\max\Bigl(\Bigl\lfloor\frac{\chi_{\mathrm{last}}\,\Delta t/\Delta t_{\mathrm{last}}}{0.04}\Bigr\rfloor+1,\
    \Bigl\lfloor\frac{\Delta t}{\Gamma\,\tau_{RC}}\Bigr\rfloor+1\Bigr),\ 400\Bigr),
  \qquad \Gamma=1,
\end{equation}
$\chi_{\mathrm{last}}$ the largest fractional vessel-pressure change over the
previous tick. The RC bound:
\begin{equation}
  C=\frac{m_u}{p_u}\ \si{\kilogram\per\pascal},\qquad
  R=\min_{\text{regulators}}\frac{D_r}{\mdot_{\mathrm{rated}}}\ \si{\pascal\second\per\kilogram},\qquad
  \tau_{RC}=\min_{\text{tanks}}\,CR .
  \label{eq:rc}
\end{equation}
Line resistance is omitted, which only lengthens $\tau_{RC}$, so the bound errs
toward more steps. $\Gamma$ was swept: tick-to-tick roughness is flat to
$\Gamma=2$ and broken by $4$; \cref{eq:rc} explains why helium (small $C$) was
forty times noisier than nitrogen at the same step. Within each coupling step
the vessels take 4 sub-steps, each further divided so that no sub-step changes
ullage mass by more than \SI{2}{\percent} or estimated temperature by more than
\SI{3}{\percent} (cap 200), with up to six halvings on an EOS refusal; the
ullage mass is floored at what atmosphere would hold in its volume.

Further rules of the live path: a tank stops accepting gas once
$p_u > p_{\mathrm{supply}}(1+0.005)$ and the refused mass is credited back to
the bottle (without this helium lost half its bottle); a tank below
\SI{1}{\gram} of liquid isolates its outlets; actuators slew linearly at
$1/t_{\mathrm{travel}}$ (default \SI{50}{\milli\second}); the warm start is
dropped for flows but kept for pressures when the isolated set changes;
indeterminate stub nodes hold their last known pressure.

\paragraph{The cockpit runs the study's numerics.} Until 2026-09-11 the
interactive path differed from the batch path in three more ways: a Newton cap
of 30 iterations against the study's 120, a single outer step per panel tick
(\SI{0.25}{\second}) against the study's \SIrange{0.01}{0.02}{\second}, and a
wall-clock budget of \SI{0.15}{\second} per tick that folded the remaining
coupling steps into one when the solver fell behind. The fold was the
difference an operator could see: traces that were smooth in a study were
rough on the console, and a burn was pre-computed at study settings and
replayed rather than run. All three are gone. A panel tick of $\Delta t$ is
integrated as $\lceil\Delta t/0.02\rceil$ outer steps of \SI{0.02}{\second}
(each with the coupling and sub-step rules above), the Newton cap is 120, and
there is no budget: a stand whose regulator--ullage loop is too stiff for real
time runs in slow motion, and the panel reports the ratio of stand time to
wall time. A pre-computed replay remains available to a client that asks for it.

\paragraph{Limits.} Each vessel carries the MAWP its drawing declares (gauge).
A step that leaves a vessel above it stops the stand on that frame: the
state is reported with the vessel and the pressure it reached, commands are
refused until the stand is reset, and the frame -- with its history -- stays
on the panel. The model could integrate on to the propellant's critical
pressure; a tank could not.

%==============================================================================
The number of coupling steps per tick is the largest of three: the
pressure-change rule ($\Delta p/p \le 0.04$ per step, from the previous tick),
the RC rule ($\Delta t \le \Gamma\,C R$), and a mass rule --- no step may move
more than \SI{4}{\percent} of any ullage's mass, sized from the last solve's
flows, because the RC rule's $R$ is the regulator's droop slope, its stiffness
near lockup, and a wide-open regulator passes ten times its rated flow into a
\SI{0.43}{\litre} ullage. Independently of step size, a vessel clips the last
step of a press onto its supply: the admitted mass is
$\Delta m \le m_u\,[(p_s(1+\beta) - p)/p]/\gamma$ with $\gamma=1.67$ bounding the
charge heating, and the remainder is refused back to the bottle. Without the
clip the fuel tank overshot lockup by \SI{35}{\psi} on the crossing step and the
regulator shut on it.

\section{Thermophysical properties}\label{sec:props}
%==============================================================================
Each species has an ordered backend chain, default (\code{bicubic}, \code{heos}):
CoolProp's tabulated bicubic interpolant over the Helmholtz EOS (10--30$\times$
faster), which \emph{refuses} outside its envelope, falling through to the EOS.
A measured $(p,T)$ table may be placed in front. Addressable pairs:
$(p,T),(p,h),(p,s),(p,q),(T,q),(p,\rho),(h,s),(\rho,T),(\rho,u)$. $Z$ is derived
as $p/(\rho R_sT)$ on every backend (agrees with CoolProp's own to $3\times10^{-6}$).
Backends are thread-confined; results are memoised on the resolved
state (256 entries, cleared wholesale). Substance constants
($p_{\mathrm{crit}},T_{\mathrm{crit}},M$) are read from the chain.

%==============================================================================
\section{From drawing to network}\label{sec:pid}
%==============================================================================
% NOTE: These rules are methodology, and several of them are where real
% defects lived. Each deserves a sentence on the failure it prevents.
\begin{enumerate}[label=(D\arabic*)]
\item \textbf{Fluid and temperature propagate together} by breadth-first walk
      from declared symbols, stopping at barriers (\code{TANK}, \code{ENGINE},
      \code{INJECTOR}); two passes --- supplies first, then tanks down their
      outlets only --- so a manifold one hop from a tank is not painted with
      propellant. A declaration always beats an inherited value.
\item \textbf{A tank is two fixed boundaries}, ullage and outlet, with no
      branch between them. The ullage node carries the \emph{pressurant}
      species. A line attaches to the ullage if its far end reaches a gas
      supply or atmosphere without crossing a vessel; otherwise the outlet.
      Every node on the pressurant or vent reach is declared \emph{gas};
      every tank outlet, injector face and node painted with a tank's contents
      is declared \emph{liquid} (\cref{eq:phaserule}).
\item \textbf{Vents are inferred:} a \code{MAN}/\code{ROT}/\code{SOL}/\code{RV}
      with exactly one unplumbed port has that port fixed at
      $p_{\mathrm{amb}}=\SI{101325}{\pascal}$.
\item \textbf{A dome-control regulator is lifted out of the flow graph}; its
      outlet setpoint, evaluated at its own supply pressure, becomes the
      \code{dome} signal of the regulator it loads --- which is how the
      supply-pressure effect reaches every tank through the control regulator's
      own $S$.
\item \textbf{The engine attaches by species:} each line reaching the engine
      symbol is re-terminated on an injector face node for the propellant it
      carries; an unmatched species is warned, not dropped.
\item \textbf{Instruments} read their clip's upstream node by default, or the
      downstream node on request; a transducer clipped to nothing reads nothing.
\item \textbf{Defaults are warned:} a source with no pressure is \SI{500}{\psi},
      a sink \SI{350}{\psi}; a tank whose fluid is gaseous at its stated
      temperature is flagged with the saturation temperature it is missing.
\end{enumerate}
Uncoupled steady builds give the tank outlet the static boundary
$p_{\mathrm{ullage}} + \rho_\ell g\,\ell$, with $\ell$ the level of a
\SI{95}{\percent} fill in the drawn volume at the drawn bore with 2:1 heads,
inverted exactly as the cockpit inverts it; a missing volume or bore is
defaulted (\SI{17.5}{\litre}, \SI{152}{\milli\metre}) and warned. The cockpit
overwrites the value with \cref{eq:outlet} at every step.

%==============================================================================
\section{Verification and validation}\label{sec:vv}
%==============================================================================

\subsection{The governing rule}
\begin{quote}
A model cannot be verified with itself. A run that converges, plots smoothly
and violates no assertion is not evidence.
\end{quote}

\subsection{Tiered regimen}
% EXPAND: one paragraph per tier; the doc is docs/PHYSICS-BENCHMARK.md.
\begin{description}[leftmargin=2em,style=nextline]
\item[Tier 0 --- gates] Type, format, unit tests.
\item[Tier 1 --- closed form] Component checks against hand calculation and the
  reference libraries, no network, no integrator: Darcy--Weisbach by hand; a
  one-segment run equals a plain pipe; choking against \code{fluids}' own IEC
  60534 (\SI{0.00000}{\percent} across five cases); latent heats against the
  handbook; boil-off exactly $q/h_{fg}$; vapour partial pressure positive and
  monotone over seven decades; chilldown magnitudes; stainless $c_w$ against
  NIST; $\mathrm{Nu}$ against \code{ht}; Biot bounds; the icicle arithmetic.
\item[Tier 2 --- integration] The He/N$_2$ study at $\Delta t=\SI{10}{\milli\second}$
  (the ignition dip is a \SI{40}{\milli\second} feature), all pressures psig.
  Expectations live in \code{docs/PHYSICS-BENCHMARK.md} \S2.1 and are
  re-baselined there whenever a physics change moves them; tolerance
  $\pm\SI{2}{\psi}$, $\pm\SI{0.05}{\second}$. Both gases start at
  \SI{550.0}{\psig}; helium recovers lockup within half a second and nitrogen
  in about three. Full-physics variant: thermal effects move the
  \emph{pad state}, not the burn (N$_2$ cases within \SI{0.4}{\psi} of each other
  at the recorded instants). Line-wall variant and its sensitivity sweep.
  % NOTE: re-baselined 2026-09-10; the previous table predated enthalpy
  % propagation and had never been re-run. Say that, and say the rule cuts
  % both ways.
\item[Tier 3 --- venting and state machine] Blowdown against a hand-integrated
  choked isentrope (\SI{30}{\milli\second} simulated vs \SI{39}{\milli\second}
  hand; the simulator must never be \emph{faster} than the sonic limit);
  isolation; vents must exist; the state machine fails closed.
\item[Tier 4 --- the trap list] Named failures, each with the check that
  catches it: reference-state mixing, invented geometry wearing a real
  provenance tag, dataclass fields dropped on rebuild, non-physical states
  reported as success, a ceiling of $0$ meaning shut, reverse flow through a
  regulator, documents destroyed by a successful-looking save, benchmarks that
  measure their own cache.
\end{description}

\subsection{Sabotage testing}
Every physics test is verified by breaking the mechanism it guards and
confirming the test fails.

\subsection{Energy audits}
Wall heat delivered to the stream, integrated, against $\int m_wc_w\,\mathrm{d}T_w$:
closure to \SI{0.45}{\percent}, the residual being midpoint-evaluated $c_w$.

\subsection{Toggle discipline}\label{sec:toggle}
\emph{Off} is byte-identical to the baseline; \emph{on} against a geometry
declaring nothing for the model is also byte-identical. The second is not
implied by the first: it requires gating on the model having something to do,
not merely on the flag. (The line-wall model reads the walked arrival enthalpy
into \cref{eq:ullage} only when \emph{both} the flag is set \emph{and} the
drawing declared metal.)

\subsection{Study methodology}
% NOTE: Say why the study primes T-0 directly instead of flying the pad
% sequence: pressing draws on the bottle under test, so an undersized COPV
% would fail for two reasons at once.
Tanks primed at \SI{95}{\percent} fill and \SI{550}{\psig} (dome \SI{500}{\psig}
$+$ \SI{50}{\psi} bias), bottle at \SI{4500}{\psig}, a stated pad hold of
\SI{300}{\second} (interface contact time; wetted wall at liquid temperature,
ullage-side wall at gas temperature); then settled with the press solenoids
held open until every tank has sat within \SI{4}{\psi} of lockup (wider than the regulator's own \SI{0.5}{\percent} dead band) for
\SI{0.5}{\second} (at least \SI{2}{\second}, at most \SI{10}{\second}, and a
note if it never does), valves released into \code{Ready}, \SI{0.5}{\second}
recorded lead-in as the ignition datum, burn to \SI{0.06}{\kilogram} remaining; the pressure floor is
taken after \SI{0.3}{\second}, converged samples only. Sweep volumes
$\{2,3,4,6,8\}\,\si{\litre}$ with the as-built bottle inserted at its real
volume.

%==============================================================================
\section{Stated limitations}\label{sec:limits}
%==============================================================================
% EXPAND: a paragraph each.
\begin{enumerate}[label=(L\arabic*)]
\item \textbf{Liquids priced saturated} (\cref{eq:phaserule}, \cref{rem:phase}).
      Subcooling is ignored. An \emph{undeclared} sub-critical node falls back
      to the temperature rule; the drawing reader declares every node it can
      place.
\item \textbf{No soak} (\cref{as:soak}).
\item \textbf{Lumped fitting walls} (\cref{sec:biot}): $\mathrm{Bi}\approx1$,
      over-credit \SIrange{19}{32}{\percent}, worth $\approx$\SI{3}{\psi} of
      \SI{54}{\psi} because doubling declared metal moves the answer by
      \SI{6}{\psi}; fails for thin, sparsely fitted, high-flow geometry.
\item \textbf{Vapour energy} (\cref{prop:reference}); \textbf{collapse is a lower
      bound} (\cref{eq:collapse}).
\item \textbf{Quasi-steady lines} (\cref{as:qs}), reported per branch.
\item \textbf{Two wall lumps split at a fixed area fraction} (\cref{eq:twowall}).
      As the level falls, cold wetted wall becomes ullage-side wall; the model
      does not hand metal from one lump to the other, so a long burn understates
      how much cold wall the ullage sees.
\item \textbf{Vessel walls are estimated unless declared.} A drawing may
      declare $m_w$, $c_w$, $(hA)_w$ on any vessel symbol. Undeclared, a tank
      takes \SI{0.457}{\kilogram\per\litre}, \SI{900}{\joule\per\kilogram\per\kelvin},
      $(hA)_w = 12\,(V/\SI{17.5}{\litre})^{2/3}\ \si{\watt\per\kelvin}$; a
      bottle \SI{1.36}{\kilogram\per\litre} (steel K-bottle), $500$,
      $20\,(V/\SI{4.687}{\litre})^{2/3}$; each use is reported as an assumption.
      $(hA)_w$ remains the least-known number on the stand, and the split of
      thermal effects between bottle wall and line walls depends on it.
\item \textbf{$x_T=0.7$ default} is the largest choked-gas uncertainty;
      $K_{\mathrm{turn}}=1.3$ manifold default is unmeasured.
\item \textbf{Chamber closure is not converged in the cockpit} (one $0.3$ step
      per tick) and is three fixed-point passes in the chamber itself.
\item \textbf{Temperature lags the solve by one step}; no viscous heating; no
      two-phase or flashing feed; a single pressurant species; GSE fill and
      bottle fill are commanded rates, not solved. The pressurant bottle
      is charged on the pad by default (optionally it arrives full and cold
      at the drawing's pressure, a supplier's cylinder); the charge is
      adiabatic from a bank at ambient, $T_{\mathrm{end}}\to\gamma T_{\mathrm{bank}}$
      less what the wall takes, and sags as it cools.
\item \textbf{One-dimensional, lumped}: no manifold maldistribution, injector
      internal flow, or acoustics.
\end{enumerate}

%==============================================================================
\appendix
\section{Constants}
%==============================================================================
\begin{center}
\small
\begin{tabular}{llll}
\toprule
Symbol / name & Value & Meaning & Where \\
\midrule
$\kappa$ & \SI{e-3}{\pascal\second\per\kilogram} & Jacobian slope floor & \cref{sec:regular} \\
$\mathcal{K}$ & \SI{e8}{\pascal\second\per\kilogram} & choked-row stiffness & \cref{eq:chokerow} \\
$R_{\mathrm{rev}}$ & \SI{e10}{\pascal\second\per\kilogram} & regulator reverse stiffness & \cref{sec:regulator} \\
tol & $10^{-6}/10^{-5}/10^{-4}$ & steady / transient / live & \cref{eq:norm} \\
$\delta$ & $\max(10^{-6}|\mdot|,10^{-9})$ & FD perturbation & \cref{eq:fd} \\
$\mdot_0$ & \SI{e-3}{\kilogram\per\second} & initial branch flow & --- \\
$\sigma_{\mathrm{shut}}$ & $10^{-3}$ & valve isolation travel & \cref{sec:reduce} \\
iterations & $100$ / $120$ & Newton cap, standalone / live (was 30 live) & \cref{sec:rc} \\
$\mathrm{Re}_{\mathrm{lam}}$ & $2040$ & laminar--turbulent, friction and Nusselt alike & \cref{eq:darcy} \\
$x_T$ & $0.7$ & default terminal ratio & \cref{eq:iec} \\
$K_v/C_v$ & $1.1561$ & flow-coefficient pairing & \cref{sec:choking} \\
$K_{\mathrm{turn}}$ & $1.3$ & manifold port turn & \cref{sec:manifold} \\
$K_{\mathrm{end}}$ & $0.5$ & flex-hose end fittings & --- \\
$C_d^\infty, C_d^{\min}$ & $0.6,\,0.15$ & injector defaults & \cref{eq:cd} \\
band & $0.005$ & regulator shut / tank supply clamp & \cref{sec:regulator}, \cref{sec:rc} \\
$t_{\min}$ & \SI{1}{\milli\second} & collapse contact-time floor & \cref{eq:collapse} \\
$h_{\mathrm{film}}, h_{\mathrm{nuc}}, \Delta T_{\mathrm{Leid}}, \Delta T_{\mathrm{onb}}$ & $100, 3000$ \si{\watt\per\square\metre\per\kelvin}; $40, 2$ \si{\kelvin} & wetted-wall regimes, boiling onset (cockpit) & \cref{eq:regimes,eq:wallboil} \\
$\delta_s, h_{\mathrm{mix}}$ & \SI{1}{\centi\metre}, \SI{5}{\watt\per\square\metre\per\kelvin} & surface layer, layer--bulk mixing (cockpit) & \cref{eq:strat} \\
$\rho_s$ & \SI{8000}{\kilogram\per\cubic\metre} & stainless & \cref{sec:wall} \\
$m_{\mathrm{union}},D_{\mathrm{ref}}$ & \SI{113}{\gram}, \SI{10.92}{\milli\metre} & fitting-mass anchor & \cref{sec:wall} \\
sweeps, tol, min flow & $12$, \SI{0.05}{\kelvin}, \SI{e-6}{\kilogram\per\second} & enthalpy walk & \cref{eq:enthalpy} \\
NTU cap & $40$ & effectiveness saturation & \cref{eq:ntu} \\
$\omega$ & $0.5/0.4$ & chamber relaxation, batch (cockpit is root-solved, tol \SI{0.5}{\psi}, 12 it.) & \cref{sec:engineloop} \\
secant guard & $f'<0.5$ & take secant step & \cref{sec:engineloop} \\
$\Lambda$ thresholds & $0.20,\,0.10$ & soft / barely metering & --- \\
rtol, atol & $10^{-6},10^{-9}$ & stiff integrator & --- \\
$\Gamma$ & $1$ & RC safety & \cref{eq:rc} \\
$\Delta t_{\max}$, outer step & \SI{0.25}{\second}, \SI{0.02}{\second} & panel tick cap; live outer step (no wall-clock budget) & \cref{sec:rc} \\
coupled change, mass, cap & $0.04$, $0.10$, $400$ & coupling-step rules & \cref{sec:rc} \\
sub-steps & $4$; $2\%$ mass, $3\%$ T; cap $200$; 6 halvings & vessel Euler & \cref{sec:rc} \\
$t_{\mathrm{travel}}$ & \SI{50}{\milli\second} & default actuator & \cref{sec:rc} \\
dry & \SI{1}{\gram} (live), \SI{60}{\gram} (study) & tank empty & --- \\
$p_{\mathrm{amb}}$ & \SI{101325}{\pascal} & ambient & \cref{sec:pid} \\
fill, bore, volume & $0.95$, \SI{152}{\milli\metre}, \SI{17.5}{\litre} & static-head defaults & \cref{sec:pid} \\
vessel wall defaults & see L6 & tank / bottle & \cref{sec:limits} \\
\bottomrule
\end{tabular}
\end{center}

\section{Correlation provenance}
% EXPAND: fill in editions and page references.
\begin{center}
\small
\begin{tabular}{lll}
\toprule
Quantity & Source & Adapter \\
\midrule
$f_D$ & Colebrook--White; Clamond (exact); others selectable & \code{fluids.friction} \\
$f_D^{\mathrm{curved}}$ & curved-pipe friction & \code{fluids.friction\_factor\_curved} \\
Fitting $K$ & Crane TP-410; Hooper 2K; Darby 3K & \code{fluids.fittings} \\
$C_v\to K$ & \cref{eq:cvk} & \code{fluids.fittings.Cv\_to\_K} \\
Gas sizing, $\mathcal{C}$ & IEC 60534-2-1 & \code{fluids.control\_valve} (calibrated) \\
Liquid choke $F_L$ & IEC 60534-2-1 & \code{fluids.control\_valve} \\
ISO orifice $C$ & ISO 5167 Reader--Harris/Gallagher & \code{fluids.flow\_meter} \\
Inherent characteristics & linear / equal-\% / quick-opening & \code{fluids.control\_valve} \\
$\mathrm{Nu}$ & Dittus--Boelter; laminar $3.66$ & \cref{eq:ntu}; checked vs \code{ht} \\
$c_w(T)$ & NIST cryogenic material data & piecewise linear \\
$c^*,C_F,T_c,\gamma$ & NASA CEA & EngineDesign cache \\
$C_d(\mathrm{Re},d)$ & EngineDesign's model, reproduced & \cref{eq:cd} \\
EOS & Helmholtz reference equations & CoolProp HEOS / BICUBIC \\
\bottomrule
\end{tabular}
\end{center}

\end{document}
